\section{Evaluation}
\label{sec:eval}

We separate three kinds of evidence. The launch experiment measures the physical
observation and effect times of the implementation. The two-GPU study compares
policies at selected live operating points. The discrete-time Markov sweep tests
the matched-mismatch law under the process assumed by the theorem. Unless noted,
we report the mean and a $t$-based $95\%$ interval over independent repeats.
Three questions organize the study. \textbf{(Q1)} What observation-to-effect lag
does the launch-on-demand path exhibit. \textbf{(Q2)} How do the feedback, static
and greedy policies behave under controlled live contention. \textbf{(Q3)} Does
the matched-mismatch value in Theorem~\ref{thm:limit} hold in the discrete-time
model, and does the online switch behave mechanically on hardware.

\subsection{Observation and effect times (Q1)}
\label{sec:eval-detection}
A read-only \texttt{nvidia-smi} occupancy probe takes $27$ to $33$ ms at the
median. The measured time from launching a real GPU worker until its utilization
first exceeds $40\%$ is $2.46$ s over 10 thermally guarded repeats on a GV100,
with median $2.48$ s and range $2.29$ to $2.61$ s. The first measurement is an
observation time for already running work. The second is a launch-to-visible-work
time for the scale-out path and instantiates Proposition~\ref{prop:launchbound}.
It does not establish a universal lower bound for warm workers or pipelined
execution.

One completion-gated admission epoch lasts about $1.1$ s on average, so the
launch path corresponds descriptively to about two or three decision epochs. The
epoch duration is policy dependent and variable. We therefore use controlled
integer values of $D$ in Q3 rather than treating the wall-clock conversion as an
exact live Markov lag.

\subsection{Selected live operating points (Q2)}
\label{sec:eval-contended}
We create controlled contention on a two-GPU node with $C=2$. A separate
experiment process occupies $c(t)\in\{0,1\}$ GPUs, leaving guest capacity
$a(t)=C-c(t)$. We refer to it as the competing workload rather than an
independent production tenant. The prober observes a per-GPU process census.
Square load uses a period-$2\tau$ wave. Poisson load uses memoryless on and off
arrivals with a matched duty cycle. Alongside useful completion rate and mismatch
fraction $\rho$, we report the competing workload's completed work relative to
its undisturbed baseline.

The study used randomized interleaved order and four seeds per cell. Parameters
fixed before the runs were $K=4$, $w_{\min}=1$, $\alpha=0.5$, $\beta=0.5$,
$\theta=0.15$, $\bar a=1$, a $60$ s square period, a Poisson rate matched to the
same duty cycle, an $8$ s GPU work unit, and $64$ submitted Jobs per trial. The
controller used a real-valued target and integer execution as specified in
Algorithm~\ref{alg:adaptive}. Each trial ended when the burst drained.

Table~\ref{tab:e9} reports the selected operating points. Under square load, the
feedback branch obtains $0.091$ tasks/s against $0.064$ tasks/s for the static
branch, a ratio of $1.41$. The mismatch fractions differ, $0.15$ against $0.11$,
so this is not a matched-constraint test of Theorem~\ref{thm:limit}. The greedy
policy obtains $0.101$ tasks/s at mismatch $0.64$ and leaves the competing
workload $62\%$ of its baseline. The feedback branch leaves $76\%$, while the
static branch leaves $89\%$.

Under Poisson load, feedback and static goodput are statistically
indistinguishable at the shown precision, $0.062$ and $0.064$ tasks/s. Feedback
has the larger mismatch fraction, $0.30$ against $0.22$. This pre-registered
break-even case is consistent with reduced predictability, but the live data do
not estimate $r(D)$ precisely enough to identify a point on
\eqref{eq:adv}. The mismatch metric is a model quantity. It is not presented as
a direct count of host-policy violations.

\begin{table}[t]
\centering
\caption{Selected live operating points with $C=2$ and four seeds. Entries are
means with $95\%$ intervals where shown. ``Kept'' is competing-workload
throughput relative to its undisturbed baseline. The feedback and static rows are
not matched in $\rho$.}
\label{tab:e9}
\small
\setlength{\tabcolsep}{3.5pt}
\begin{tabular}{llccc}
\toprule
contention & policy & goodput (tasks/s) & $\rho$ & kept \\
\midrule
\multirow{3}{*}{square}
 & greedy   & $0.101\pm0.002$ & $0.64\pm0.01$ & $0.62$ \\
 & feedback & $0.091\pm0.000$ & $0.15\pm0.01$ & $0.76$ \\
 & static   & $0.064\pm0.000$ & $0.11\pm0.01$ & $0.89$ \\
\midrule
\multirow{3}{*}{Poisson}
 & greedy   & $0.096\pm0.002$ & $0.85\pm0.06$ & $0.62$ \\
 & feedback & $0.062\pm0.006$ & $0.30\pm0.06$ & $0.55$ \\
 & static   & $0.064\pm0.000$ & $0.22\pm0.04$ & $0.79$ \\
\bottomrule
\end{tabular}
\end{table}

\subsection{Matched-mismatch model test and live switch check (Q3)}
\label{sec:eval-law}
We next simulate the window controllers on the symmetric two-state Markov process
used by Theorem~\ref{thm:limit}. The per-epoch flip probability $p$ sweeps the
coherence time, the lag is $D=3$, and each point contains $1.2\times10^5$ epochs.
For every controller, the plotted quantity is goodput increase over the best
record-independent policy at the same mismatch fraction.

Figure~\ref{fig:tausweep} supports four model-level readings. First, the delayed
threshold policy follows $r(D)/2$ with maximum displayed deviation $0.013$.
Second, fixed AIMD captures only part of the available increase at long coherence
time and becomes worse than a matched static policy as predictability falls.
Third, the branch-adaptive controller stays above the static branch at the
simulated points and recovers much of the AIMD increase where its estimate exceeds
the switch level. This is an empirical property of the plotted parameterization,
not a pointwise theorem. Fourth, a tabular policy-gradient agent with four delayed
readings converges within $0.003$ of the threshold policy at every plotted point
and exceeds the analytical line by at most the same $0.013$ sampling deviation.
Under the Markov model, the older readings are conditionally redundant given the
newest one. The test therefore checks whether learning recovers the sufficient
statistic.

The right panel sweeps $D$ and shows the expected collapse against $D/\tau_c$.
These are single long simulations rather than seed-based confidence intervals.
Their role is to verify the model implementation and the matched-constraint
calculation, not to substitute for a live frontier.

\begin{figure}[t]
\centering
\includegraphics[width=\columnwidth]{figures/tau_sweep.png}
\caption{Discrete-time controller study. \emph{Left.} Goodput increase over a
matched static policy at $D=3$. The delayed threshold policy follows
$r(D)/2$. Fixed AIMD loses value as predictability falls, and the branch-adaptive
controller selects between AIMD and static at the plotted points. \emph{Right.}
Sweeping $D$ collapses the threshold increase against $D/\tau_c$.}
\label{fig:tausweep}
\end{figure}

\noindent\textbf{Live behavior of the switch.}
We also ran Algorithm~\ref{alg:adaptive} on the two-GPU node across a Markov
contention sweep with $p$ from $0.01$ to $0.4$ and two seeds. The sliding-window
estimate satisfies $|\hat r-(1-2p)^D|\le0.07$ over the reported points. The
fraction of epochs using the feedback branch falls from $0.98$ to $0.48$ as
predictability crosses the switch level. This live sweep does not reproduce a
goodput crossover. Raw goodput rewards mismatch, so always-on AIMD outperforms
the static branch at every live point by accepting much higher mismatch at large
$p$, $0.44$ against $0.04$. The adaptive switch instead sheds mismatch relative
to always-on AIMD, with the reduction growing from $0.01$ to $0.22$ as $p$ rises
from $0.08$ to $0.4$. With two seeds, this is a mechanical check of the estimator
and branch actuation only.

\subsection{Summary and threats to validity}
\label{sec:eval-threats}
The measurements establish the timing of the launch-on-demand path, one useful
structured-contention operating point, one break-even operating point, and a
matched-constraint verification in the discrete-time model. They do not establish
a live value of $r(D)/2$.

The main limitation is the absence of a live matched-$\rho$ frontier. The static
policy in Table~\ref{tab:e9} is one operating point rather than the best
record-independent policy at the same mismatch. A stronger experiment would
sweep randomized static admission probabilities and feedback parameters, then
compare the maximum goodput at common values of $\rho$. Contention is controlled
on one two-GPU node rather than drawn from a production multi-user trace. Four
seeds support Table~\ref{tab:e9}, while the live switch sweep has two. The
``kept'' column has no interval, and values printed as $0.000$ are below the shown
resolution rather than exactly zero. Completion-gated epochs also make the live
wall-clock lag action dependent. These limits are why the exact theorem is tested
in the controlled discrete-time sweep rather than claimed from the hardware
numbers.
